% !TEX root = main.tex
\section{Collatz Source Value}\label{sec:collatz-source-value}
Now we are going to explore the trace back value according to \cref{sec:extend-the-basic-formula}.
\begin{definition}\label{def:source-value-function}
From \cref{theorem:any-presequence}, with the same arguments $\alpha$, $\beta$ and $\gamma$, and target value $b\in \widehat{\mathbb{Q}}_{{\beta}}[b]$, for any finite vector $\bm{v} = (v_{1}, v_{2}, \cdots, v_{n}) \in \mathbb{N}^{n}$, we can derive a sequence $\{v_i\}_{i=1}^{n}$ and its accmulation $\mathcal{A}$. With this we can also derive Collatz complement $\mathfrak{V}$ and reduced Collatz complement $\mathfrak{v}$ by arguments $(\alpha, \beta, \gamma)$. And then we have a function $f_n : \mathbb{N}^{n} \to \mathbb{Q}$, that defined as
\[ f_{n}(\bm{v};b) = b \cdot \frac{\beta^{\mathcal{V}_{n}}}{\alpha^{n}} - \frac{\gamma}{\alpha}\cdot \mathfrak{v}_{n} \in \widehat{\mathbb{Q}}_{{\beta}}[b] \]
is called the source value from $b$ by $\bm{v}$. In the absence of ambiguity, we denote $f_{n}(\bm{v}) = f_{n}(\bm{v};b)$.
\end{definition}

And to describe the iterative properties of Collatz sequences, we denote $a \overset{k}{\rightsquigarrow} b$ if there exists $k$ such that $a_k = b$.
If it $k$ doesn't matter, we denote $a \rightsquigarrow b$. Notice that we always suppose the step count of $a \rightsquigarrow b$ will never be $0$. If we need to represent $a_0 = b$, we will denote it precisely as $a \overset{0}{\rightsquigarrow} b$. We say $a \not\rightsquigarrow b$, means there is not $n > 0$, such that $a \overset{n}{\rightsquigarrow} b$.
With these we have
\begin{proposition}
    $f_{n}(\bm{v}) \overset{n}{\rightsquigarrow} b$.
\end{proposition}
\begin{proposition}
    $a\rightsquigarrow a$ and $a \rightsquigarrow b$, then $b \rightsquigarrow a \rightsquigarrow b$.
\end{proposition}
\begin{proposition}
    If $\mathfrak{v}_{n}$ follow the requirements of \cref{theorem:any-presequence}, then we will have $\beta \nmid f_{n}(\bm{v})$.
\end{proposition}
\begin{proposition}
    $f_{n}(\bm{v}) = f_{n}(\bm{u}) \Longleftrightarrow \bm{v} = \bm{u}$.
\end{proposition}
\begin{proposition}
    $f_{n}(\bm{v}) = \frac{\beta^{v_{1}}}{\alpha}\left(\frac{\beta^{v_{2}}}{\alpha}\left(\cdots \left(\frac{\beta^{v_{n}}}{\alpha}\cdot b - \frac{\gamma}{\alpha}\right)\cdots\right)- \frac{\gamma}{\alpha}\right) - \frac{\gamma}{\alpha}$.
\end{proposition}
\begin{proposition}\label{prop:condition-of-source-value-le-N}
    Given $N > 0$, $f_{n}(\bm{v}) \le N \Longleftrightarrow \frac{\alpha^{n}\cdot N + \gamma\cdot\mathfrak{V}_{n}}{a\cdot \beta^{\mathcal{V}_{n}}} \ge 1$.
\end{proposition}
We denote $Z_{n}(N;\bm{v}) = \frac{\alpha^{n}\cdot N + \gamma\cdot\mathfrak{V}_{n}}{a\cdot \beta^{\mathcal{V}_{n}}}$, and we will see, it will play very import role in \cref{sec:distribution-of-source-value-set}.
\begin{theorem}\label{thm:vary-rule}
Given $\Delta\bm{v} \in \mathbb{Z}^{n}$, view it as a $n$-length Collatz reduce sequence. Then we have
\[ f_{n}(\bm{v} + \Delta\bm{v}) = \beta^{\Delta\mathcal{V}_{n}}\cdot f_{n}(\bm{v}) + \frac{\gamma}{\alpha^{n}} \cdot \nabla\cdot\mathfrak{V}_{n} = \beta^{\Delta\mathcal{V}_{n}}\cdot f_{n}(\bm{v}) + \frac{\gamma}{\alpha} \cdot \nabla\cdot\mathfrak{v}_{n}. \]
Here
\begin{align*}
    \nabla\cdot\mathfrak{V}_{n} &= \sum_{k=0}^{n-1}\beta^{\mathcal{V}_{k} + \Delta\mathcal{V}_{k}}\cdot\alpha^{n-k-1}\cdot(\beta^{\Delta\mathcal{V}_{n}^{k}} - 1),\\
    \nabla\cdot\mathfrak{v}_{n} &= \sum_{k=0}^{n-1}\beta^{\mathcal{V}_{k} + \Delta\mathcal{V}_{k}}\cdot\frac{\beta^{\Delta\mathcal{V}_{n}^{k}} - 1}{\alpha^{k}}.
\end{align*}
\end{theorem}
\begin{proof}
    Let $\bm{u} = \bm{v} + \Delta\bm{v}$. First we have
    \begin{align*}
        \mathfrak{U}_{n} &= \beta^{\Delta\mathcal{V}_n}\cdot\mathfrak{V}_n - \nabla\cdot\mathfrak{V}_{n},\\
        \mathfrak{u}_{n} &= \beta^{\Delta\mathcal{V}_n}\cdot\mathfrak{v}_n - \nabla\cdot\mathfrak{v}_{n}.
    \end{align*}
    Notice that
    \[f_{n}(\bm{u}) = b \cdot \frac{\beta^{\mathcal{U}_{n}}}{\alpha^{n}} - \frac{\gamma}{\alpha}\cdot \mathfrak{u}_{n} = \beta^{\Delta\mathcal{V}_n} \cdot b\cdot \frac{\beta^{\mathcal{V}_{n}}}{\alpha^{n}} - \frac{\gamma}{\alpha}\cdot \mathfrak{u}_{n},\]
    and
    \[ \mathfrak{u}_{n} = \sum_{k=0}^{n-1}\frac{\beta^{\mathcal{U}_k}}{\alpha^{k}} = \sum_{k=0}^{n-1}\frac{\beta^{\mathcal{V}_k + \Delta\mathcal{V}_k}}{\alpha^{k}}. \]
    Then
    \begin{align*}
        & f_{n}(\bm{u}) - \beta^{\Delta\mathcal{V}_{n}}\cdot f_{n}(\bm{v}) = \frac{\gamma}{\alpha}\cdot\left(\beta^{\Delta\mathcal{V}_{n}}\cdot\mathfrak{v}_{n} - \mathfrak{u}_{n} \right) \\
        =&\frac{\gamma}{\alpha^{n}}\cdot\left(\beta^{\Delta\mathcal{V}_{n}}\cdot\sum_{k=0}^{n-1}\beta^{\mathcal{V}_{k}}\cdot\alpha^{n-k-1} - \sum_{k=0}^{n-1}\beta^{\mathcal{V}_{k}+\Delta\mathcal{V}_{k}}\cdot\alpha^{n-k-1} \right) \\
        =&\frac{\gamma}{\alpha^{n}}\cdot\sum_{k=0}^{n-1}\beta^{\mathcal{V}_{k}}\cdot(\beta^{\Delta\mathcal{V}_{n}} - \beta^{\Delta\mathcal{V}_{k}})\cdot\alpha^{n-k-1} \\
        =&\frac{\gamma}{\alpha^{n}}\cdot\sum_{k=0}^{n-1}\beta^{\mathcal{V}_{k}}\cdot\alpha^{n-k-1}\cdot\beta^{\Delta\mathcal{V}_{k}}\cdot(\beta^{\Delta\mathcal{V}_{n}^{k}} - 1) \\
        =&\frac{\gamma}{\alpha^{n}}\cdot\sum_{k=0}^{n-1}\beta^{\mathcal{V}_{k} + \Delta\mathcal{V}_{k}}\cdot\alpha^{n-k-1}\cdot(\beta^{\Delta\mathcal{V}_{n}^{k}} - 1)\\
        =&\frac{\gamma}{\alpha^{n}}\cdot \nabla\cdot\mathfrak{V}_{n}
        \\
        =&\frac{\gamma}{\alpha}\cdot \nabla\cdot\mathfrak{v}_{n}.
    \end{align*}
    Finally, this proposition has been proved.
\end{proof}
\begin{remark}
Though $\Delta\bm{v} \in \mathbb{Z}^{n}$, we still require $\bm{v} + \Delta\bm{v} \in\mathbb{N}^{n}$.
\end{remark}
\begin{corollary}
    Suppose $f_{n}(\bm{v})\in\mathbb{N}$, then $f_{n}(\bm{v} + \Delta\bm{v})\in\mathbb{N} \Leftrightarrow \alpha^{n} \mid \gamma\cdot\nabla\cdot\mathfrak{V}_{n}$.
\end{corollary}

\cref{thm:vary-rule} offers a very important insight.
That is, when $f_{n}(\bm{v})\overset{n}{\rightsquigarrow} b$ and $f_{n}(\bm{v} + \Delta\bm{v}) \overset{n}{\rightsquigarrow} b$, starting from a certain bit $m$ and $m + \Delta\mathcal{V}_n$, the $\beta$-carrying terms of these two numbers satisfy a certain relationship. Mathematically, that's
\begin{corollary}\label{cor:high-order-bits-sync}
    Suppose $f_{n}(\bm{v}) = \sum_{k=s}^{\infty}x_k\beta^{k}$ and $f_{n}(\bm{v} + \Delta\bm{v}) = \sum_{k=s}^{\infty}x'_k\beta^{k}$, $s\in\mathbb{Z}$. If $\alpha^{n} \mid \gamma\cdot\nabla\cdot\mathfrak{V}_{n}$, then $\exists m \in \mathbb{Z}: m \ge s$, such that
    \[ \sum_{k=0}^{\infty}x'_{k + m}\beta^{k} = \sum_{k=0}^{\infty}x_{k + m - \Delta\mathcal{V}_n}\beta^{k} + \delta, \]
    where the $\delta\in\{0, 1\}$. Vice versa.
\end{corollary}
\begin{proof}
    When $\nabla\cdot\mathfrak{V}_{n} = 0$, it's trivial. Consider $\nabla\cdot\mathfrak{V}_{n} \ne 0$.\newline
    Since $\alpha^{n} \mid \gamma\cdot\nabla\cdot\mathfrak{V}_{n}$, then $\beta^{\vert\Delta\mathcal{V}_n\vert}\cdot\frac{\gamma}{\alpha}\cdot\nabla\cdot\mathfrak{v}_{n} \in \mathbb{N}$, that means $\deg_{\beta}(\frac{\gamma}{\alpha}\cdot\nabla\cdot\mathfrak{v}_{n})<\infty$. Now let $m = \max\{\deg_{\beta}(\frac{\gamma}{\alpha}\cdot\nabla\cdot\mathfrak{v}_{n}), \Delta\mathcal{V}_n\} + s + 1$, then
    \[ \beta^{m - 1}\le\beta^{\Delta\mathcal{V}_n}\sum_{k=s}^{m - \Delta\mathcal{V}_n - 1}x_k\beta^{k} + \frac{\gamma}{\alpha}\cdot\nabla\cdot\mathfrak{v}_{n} \le \beta^{m}. \]
    And from \cref{thm:vary-rule}, we have
    \[ \sum_{k=s}^{\infty}x'_k\beta^{k} = \beta^{\Delta\mathcal{V}_n}\sum_{k=s}^{m - \Delta\mathcal{V}_n - 1}x_k\beta^{k} + \frac{\gamma}{\alpha}\cdot\nabla\cdot\mathfrak{v}_{n} + \beta^{\Delta\mathcal{V}_n}\sum_{k=m}^{\infty}x_{k}\beta^{k}. \]
    Let's right shift both the series by $m$ then we have
    \[ \sum_{k=m}^{\infty}x'_k\beta^{k - m} = \sum_{k=m - \Delta\mathcal{V}_n}^{\infty}x_{k}\beta^{k-m + \Delta\mathcal{V}_{n}} + \delta. \]
    Adjusting the subscript yields
    \[ \sum_{k=0}^{\infty}x'_{k + m}\beta^{k} = \sum_{k=0}^{\infty}x_{k + m - \Delta\mathcal{V}_n}\beta^{k} + \delta. \]
    On the other hand, if there is $m$ make the series equation holds, then $\forall k \ge m$ the series equation also holds. So suppose $m \ge \Delta\mathcal{V}_n + s + 1$. We can see from the previous derivation that
    \[\frac{\gamma}{\alpha}\cdot\nabla\cdot\mathfrak{v}_{n} = \sum_{k=s}^{m-1}x'_k\beta^{k} - \beta^{\Delta\mathcal{V}_n}\sum_{k=s}^{m - \Delta\mathcal{V}_n - 1}x_k\beta^{k} + \delta\cdot\beta^{m}, \]
    thus $\deg_{\beta}(\frac{\gamma}{\alpha}\cdot\nabla\cdot\mathfrak{v}_{n}) \le m \le\infty$. Notice that $\gcd(\alpha, \beta) = 1$, thus $\alpha^{n}\mid \gamma\cdot\nabla\cdot\mathfrak{V}_{n}$.
\end{proof}
Notice that we don't require the $f_{n}(\bm{v})$ and $f_{n}(\bm{v} + \Delta\bm{v})$ in $\mathbb{N}$, and we don't require $m \ge 0$. that means they can be any $\beta$-adic integers or $\beta$-adic finite decimals.
\begin{remark}
    If $f_{n}(\bm{v}) \not\in \mathbb{Z}_{<0}$, then $\exists m > \max\{\deg_{\beta}(\frac{\gamma}{\alpha}\cdot\nabla\cdot\mathfrak{v}_{n}), \Delta\mathcal{V}_n\} + s + 1$, such that \cref{cor:high-order-bits-sync} holds and $\delta = 0$. But if $f_{n}(\bm{v}) \in \mathbb{Z}_{<0}$ and there is
    \[ \frac{\gamma}{\alpha}\cdot\nabla\cdot\mathfrak{v}_{n} \ge \left\vert \beta^{\Delta\mathcal{V}_n}\cdot f_{n}(\bm{v}) \right\vert, \]
    then $\forall m$ large enough, there is only $\delta = 1$, i.e. $f_{n}(\bm{v} + \Delta\bm{v}) \in \mathbb{N}$.
\end{remark}
Actually, we can derive this theorem directly from \cref{theorem:any-presequence}.
Since there is $\gcd(\alpha, \beta) = 1$, the $\alpha^{-n}$ is a rational number. $\alpha^{n} \mid \gamma\cdot\nabla\cdot\mathfrak{V}_{n}$ means $\frac{\gamma}{\alpha}\cdot\nabla\cdot\mathfrak{v}_{n} \in \mathbb{N}$ and will not change the fraction part of the mixed fraction expression of the $\beta^{\Delta\mathcal{V}_{n}}\cdot f_{n}(\bm{v}) + \frac{\gamma}{\alpha} \cdot \nabla\cdot\mathfrak{v}_{n}$.
In $\beta$-adic integer $\mathbb{Z}_{\beta}$, it means the higher bits will not changed.
And the $\beta^{\Delta\mathcal{V}_{n}}$ determinate the shift count, it also will not change the highter bits. That's exactly what the theorem tell us.

Even though, we cannot determine the source value of the integer solely based on $\nabla\cdot\mathfrak{V}_{n}\in\mathbb{N}$ without the requirement of $f_{n}(\bm{v})\in\mathbb{N}$. And Even though we have $f_{n}(\bm{v})\in\mathbb{N}$, the demonstrates that solving the problem solely through analysis of $\nabla\cdot\mathfrak{V}$ is too difficult. Therefore, we must find another way to work around it.

Now let's revisit \cref{thm:vary-rule}. Analyzing the full $\Delta\bm{v}$ is too complex. If there exists a vector $\Delta\bm{v}$ with only one non-zero component, i.e., $\Delta{\bm{v}} = s\bm{e}_{k}$, such that $\bm{v} + \Delta\bm{v} = \bm{v}+s\bm{e}_{k}$, where $\bm{e}_k$ is the $k$-th $n$-dim unit vector, then the analysis becomes much simpler. With this we have
\begin{corollary}\label{cor:single-vary-rule}
$f_{n}(\bm{v}+s\bm{e}_{k}) = \beta^{s}\cdot f_{n}(\bm{v}) + \frac{\gamma}{\alpha}\cdot\mathfrak{v}_{k}\cdot(\beta^{s} - 1)$.
\end{corollary}
\begin{corollary}
$\frac{\gamma}{\alpha}\cdot\mathfrak{v}_{k} = \frac{f_{n}(\bm{v}+s\bm{e}_{k}) - \beta^{s}\cdot f_{n}(\bm{v})}{\beta^{s} - 1}$.
\end{corollary}
\begin{corollary}
    If $f_{n}(\bm{v}) = \frac{q}{\alpha^{m}}$, then $f_{n}(\bm{v}+s\bm{e}_{k}) \in \mathbb{N}$ only when $s$ holds
    \[ \beta^{s}\cdot q\cdot \alpha^{k} + (\beta^{s} - 1) \cdot \mathfrak{V}_{k}\cdot \alpha^{m} \equiv 0. \pmod{\alpha^{m+k}} \]
\end{corollary}
\begin{theorem}\label{thm:the-periodic-rule-of-reduce-sequence}
    If $f_{n}(\bm{v}) \in \mathbb{N}$, then we have
    \[f_{n}(\bm{v}+s\bm{e}_{k}) \in \mathbb{N}\Longleftrightarrow \Phi(\alpha,\beta,\gamma; k) \mid s. \]
    Here the basic adjust function is defined as
    \[ \Phi(\alpha, \beta, \gamma; k) = \alpha^{k}\cdot \widetilde{\rord_{\alpha}\beta}\cdot\Psi(\alpha,\beta,\gamma). \]
    The tilde radical order item is $\widetilde{\rord_{\alpha}\beta} = \rord_{\alpha}\beta$ if $\alpha^{k} \nmid \gamma$ else $1$.
    The auxiliary function $\Psi_{\alpha}(\beta)$ is defined as
    \[ \Psi(\alpha,\beta,\gamma) = \prod_{p\in\mathbb{P}}^{p \mid \alpha}\frac{1}{\pow(p\mid\gamma\cdot2^{\delta})\cdot\pow(p\mid\beta^{\ord_{p}\beta} - 1)},
    \]
    and $\delta = \val_2(\beta + 1) - 1$ if $2 \mid \alpha$ else $0$ .
\end{theorem}
\begin{proof}
    From \cref{cor:single-vary-rule}, we have
    \[ f_{n}(\bm{v}+s\bm{e}_{k}) = \beta^{s}\cdot f_{n}(\bm{v}) + \frac{\gamma}{\alpha}\cdot\mathfrak{v}_{k}\cdot(\beta^{s} - 1). \]
    Since $\beta^{s}\cdot f_{n}(\bm{v}) \in \mathbb{N}$, then $\frac{\gamma}{\alpha}\cdot\mathfrak{v}_{k}\cdot(\beta^{s} - 1) \in \mathbb{N}$, that equivalents
    \[\alpha^{k} \mid \gamma\cdot\mathfrak{V}_{k}\cdot(\beta^{s} - 1).\]
    Since $\gcd(\alpha, \mathfrak{V}_{k}) = 1$, the divisibility condition simplifies to
    \[\alpha^{k} \mid \gamma\cdot(\beta^{s} - 1).\]
    When $\alpha^{k} \nmid \gamma$.
    By the fundamental theorem of arithmetic, we can derive that $\alpha = \prod_{p \in \Omega(\mathbb{P} \mid \alpha)} \raf_{p}\alpha$. Thus, the global condition is equivalent to a set of local conditions for each prime $p \mid \alpha$ we have
    \[ (\raf_{p}\alpha)^{k} \mid \gamma\cdot(\beta^{s} - 1). \]
    Notice that $p \mid (\beta^{\ord_{p}\beta} - 1)$, apply the Lifting The Exponent Lemma, consequently, the inequality becomes
    \[ k\cdot\val_{p}\alpha \le \val_{p}\gamma + \val_{p}\left(\frac{s}{\ord_{p}\beta}\right) + \val_{p}(\beta^{\ord_{p}\beta} - 1) + \delta_2(\beta). \]
    Here $\delta_2(\beta) = \val_2(\beta + 1) - 1$ if $p = 2$ else $0$.
    Rearranging for $\val_{p}\left(\frac{s}{\ord_{p}\beta}\right)$, we obtain
    \[ \val_{p}\left(\frac{s}{\ord_{p}\beta}\right) \ge \val_{p}(\alpha^{k}) - \val_{p}\gamma - \val_{p}(\beta^{\ord_{p}\beta} - 1) - \delta_2(\beta). \]
    Isolating the valuation of $s$, we obtain
    \[ \val_{p}s \ge \val_{p}\left(\frac{\alpha^{k}\cdot\ord_{p}\beta}{\gamma\cdot(\beta^{\ord_{p}\beta} - 1)\cdot2^{\delta}}\right), \]
    which equivalents
    \[ \left.\pow\left(p\mid\frac{\alpha^{k}\cdot\ord_{p}\beta}{\gamma\cdot(\beta^{\ord_{p}\beta} - 1)\cdot2^{\delta}}\right)\ \right\vert\ s. \]
    From \cref{prop:raf-completely-multiplicative}, we have
    \[ \pow\left(p\mid\frac{\alpha^{k}\cdot\ord_{p}\beta}{\gamma\cdot(\beta^{\ord_{p}\beta} - 1)\cdot2^{\delta}}\right) = \frac{\pow(p\mid\alpha)^{k}}{\pow(p\mid\gamma\cdot2^{\delta})}\cdot\frac{\pow(p\mid\ord_{p}\beta)}{\pow(p\mid\beta^{\ord_{p}\beta} - 1)}. \]
    Notice that $\pow(p\mid\ord_{p}\beta) = 1$， thus we have
    \[ \left. \frac{\pow(p\mid\alpha)^{k}}{\pow(p\mid\gamma\cdot2^{\delta})\cdot\pow(p\mid\beta^{\ord_{p}\beta} - 1)}\ \right\vert\ s. \]
    Besides, we've assumed $\ord_{p}\beta \mid s$. Aggregate these all $p$ conditions, it follows
    \[ \left.\lcm_{p\in\mathbb{P}}^{p\mid \alpha}\{\ord_{p}\beta\} = \rord_{\alpha}\beta\ \right\vert\ s. \]
    When $\alpha^{k}\mid \gamma$, of course we do not need $\ord_{p}(\beta) \mid s$, thus not need $\rord_{\alpha}\beta\mid s$.
    So we let
    \[ \widetilde{\rord_{\alpha}\beta} = \left\{\begin{aligned}
        &\rord_{\alpha}\beta, && \alpha^{k}\nmid \gamma, \\
        &1, && \alpha^{k}\mid \gamma,
    \end{aligned}\right. \]
    which contains both conditions of $\alpha^{k}\mid \gamma$ or $\alpha^{k}\nmid \gamma$.
    Finally, sum up, we have $\Phi(\alpha,\beta,\gamma;k) \mid s$.
\end{proof}
In the absence of ambiguity, we denote $\Phi(k) = \Phi(\alpha, \beta, \gamma; k)$.
We have to mention that the divisibility in this proof is not the common divisibility between integers, instead of the \cref{def:extend-divisibility}. It implies that, although the derived value of $\Phi(k)$ may be a fraction, we require $s \in \mathbb{N}$, and $s$ must be a multiple of the numerator of that fractional value in its reduced form. Since the denumerator of all $\Phi(k)$ are $\Psi(\alpha,\beta,\gamma)$ and we will not pay too much attention on it.
Thus we denote $\Phi_k$ as the numerator of $\Phi(k)$'s reduced form.
From the proof, we can deduce that
\begin{corollary}\label{cor:alpha-k-mid-adjustment}
    $\forall s \in\mathbb{N}$, we have $\alpha^{k} \mid \gamma\cdot(\beta^{s\cdot\Phi_k} - 1)$.
\end{corollary}
\begin{corollary}\label{cor:alpha-times-between-Phis}
    $\forall i \ge j \ge 0$, we have $\Phi_j \mid \Phi_i$ and $\Phi_i \mid \alpha^{i-j} \Phi_j$.
\end{corollary}
\begin{example}\label{exm:sequence-of-Phi}
    Suppose $\alpha\in\mathbb{P}$ and $\val_{\alpha}(\gamma) = k$, then we have
    \[ \{\Phi_i\}_{i=1}^{n} = \{\underbrace{1, 1, \cdots, 1}_{k \text{ count}}, \Phi_{k+1}, \alpha\Phi_{k+1}, \alpha^{2}\Phi_{k+1}, \cdots, \alpha^{n-k-1}\Phi_{k+1}\}. \]
\end{example}
\begin{example}
    Suppose $q \in\widehat{\mathbb{N}}_{2}$ that $q > 1$, then we have \begin{align*}
        &\Psi(q, 2, 1) = \prod_{p\in\Omega(\mathbb{P}\mid q)}\frac{1}{\raf_{p}(2^{\ord_{p}2} - 1)}, \\
        &\Phi(q, 2, 1; k) = q^{k}\cdot\prod_{p\in\Omega(\mathbb{P}\mid q)}\frac{1}{\raf_{p}(2^{\ord_{p}2} - 1)} \cdot \lcm_{p\in\mathbb{P}}^{p\mid q}\{\ord_{p}2\}.
    \end{align*}
    They correspond to the case of $qn+1$ reduced by $2$ explored in Santos's paper \cite{santos2021collatzgeneralproblemqn1}.
\end{example}
\begin{example}
    Suppose $p, q \in \mathbb{P}$ that $p\ne q$ and $p\ne 2$, then we have \begin{align*}
        &\Psi(p, q, 1) = \frac{1}{\raf_{p}(q^{\ord_{p}q} - 1)}, \\
        &\Phi(p, q, 1; k) = \ord_{p}q \cdot \frac{p^{k}}{\raf_{p}(q^{\ord_{p}q} - 1)}.
    \end{align*}
    They correspond to the case of $pn+1$ reduced by $q$. We can observe that $\Phi_k$ must take the form $\ord_{p}q \cdot p^{m}$, where $m < k$.
\end{example}
\begin{example}
    $\Psi(3,2,1) = \frac{1}{3}$, and $\Phi(3, 2, 1; k) = 2 \cdot 3^{k - 1}$. They correspond to the classic Collatz problem.
\end{example}
\begin{example}
    $\Psi(5,2,1) = \frac{1}{5}$, and $\Phi(5, 2, 1; k) = 4 \cdot 5^{k - 1}$.
\end{example}
\begin{corollary}\label{cor:remainder-after-change-reduce-vector-modulo-alpha}
    Suppose $\gcd(\alpha,\gamma) = 1$.
    If $f_{n}(\bm{v}) \in \mathbb{N}$, then $\forall s\in \mathbb{N}$, we have
    \[ f_{n}(\bm{v}+s\Phi_{k}\bm{e}_k) \equiv f_{n}(\bm{v}) +  \frac{\gamma}{\alpha^{k}}\cdot\left(\beta^{s\Phi_k} - 1\right). \pmod{\alpha} \]
\end{corollary}
\begin{proof}
    Since $\gcd(\alpha, \gamma) = 1$, then $\alpha\mid(\beta^{s\Phi_k} - 1)$. And we  have $\mathfrak{V}_{k} \equiv 1 \pmod{\alpha}$, thus from \cref{cor:single-vary-rule}, we have
    \begin{align*}
        f_{n}(\bm{v}+s\Phi_{k}\bm{e}_k) &= \beta^{s\Phi_k}\cdot f_{n}(\bm{v}) + \frac{\gamma}{\alpha}\cdot\mathfrak{v}_{k}\cdot(\beta^{s\Phi_k} - 1) &&\\
        &= \beta^{s\Phi_k}\cdot f_{n}(\bm{v}) + \mathfrak{V}_{k}\cdot\frac{\gamma}{\alpha^{k}}\cdot(\beta^{s\Phi_k} - 1) &&\\ 
        &\equiv f_{n}(\bm{v}) +  \frac{\gamma}{\alpha^{k}}\cdot\left(\beta^{s\Phi_k} - 1\right). &&\pmod{\alpha}
    \end{align*}
\end{proof}

From \cref{cor:single-vary-rule} and \cref{cor:alpha-k-mid-adjustment}, we can observe that when $\Delta\bm{v} = \Phi_k\bm{e}_k$, the $\nabla\cdot\mathfrak{V}_{k}$ and $\nabla\cdot\mathfrak{v}_{k}$ depend on $\frac{1}{\alpha^{k}}\cdot(\beta^{\Phi_k} - 1)$.
\begin{definition}\label{def:basic-fill-part-of-adjustment}
we denote $\nabla_{k} = \frac{\beta^{\Phi_k} - 1}{\alpha^{k}}$ called the basic fill part of adjustment.
\end{definition}
Similar with Collatz complements, here we don't consider put $\gamma$ into $\nabla_{k}$. Actually, when $k = 1$, it's a kind of broad Fermat quotient, hence, in some sense, we can named them Collatz quotients, too.
\begin{definition}
If $\alpha\in\mathbb{P}_{\ge 2}$, $\alpha\nmid\gamma$ and $\alpha\nmid \nabla_{1}$, then we say the Collatz system is simple or primitive.
\end{definition}
We can verify that classic Collatz problem $(\alpha, \beta, \gamma) = (3,2,1)$ is a simple Collatz system. And simple Collatz systems will be the main objects to analyze.
\begin{proposition}
    For simple Collatz systems, we have \begin{align*}
        &\Psi(\alpha, \beta, \gamma) = \frac{1}{\raf_{\alpha}(\beta^{\ord_{\alpha}\beta} - 1)} = \frac{1}{\alpha}, \\
        &\Phi(\alpha, \beta, \gamma; k) = \ord_{\alpha}\beta\cdot\Psi(\alpha, \beta, \gamma)\cdot\alpha^{k} = \ord_{\alpha}\beta\cdot\alpha^{k-1}.
    \end{align*}
    And thus $\Phi_k = \ord_{\alpha}\beta\cdot\alpha^{k-1}$.
\end{proposition}
\begin{proposition}\label{prop:nabla-are-equivalent-modulo-alpha}
    For simple Collatz systems, $\nabla_{k} \equiv \nabla_{1} \pmod{\alpha}$.
\end{proposition}
\begin{proof}
    From \cref{cor:alpha-k-mid-adjustment} we can derive that $\alpha^{k} \mid \gamma\cdot (\beta^{\Phi_k} - 1)$. Since in simple Collatz systems, there is $\alpha \nmid \gamma$, thus $\alpha^{k} \mid (\beta^{\Phi_k} - 1)$, then $\nabla_{k} \in \mathbb{N}$. And we have $\beta^{\Phi_1} = 1 + \alpha\nabla_{1}$. And then $\nabla_{k} = \frac{1}{\alpha^{k}}((\beta^{\Phi_1})^{\alpha^{k-1}} - 1) = \frac{1}{\alpha^{k}}((1 + \alpha\nabla_1)^{\alpha^{k-1}} - 1)$. Thus we have
    \begin{align*}
        \nabla_{k} = \nabla_1 + \frac{1}{\alpha^{k}}\sum_{i=2}^{\alpha^{k-1}}\binom{\alpha^{k-1}}{i} \alpha^{i}\nabla_1^{i} \equiv \nabla_1. \pmod{\alpha}
    \end{align*}
\end{proof}
\begin{theorem}\label{thm:simple-collatz-system-mod-eq-plus-s-gamma-nabla_1}
    For simple Collatz system, we have
    \[ f_{n}(\bm{a} + s\Phi_k\bm{e}_k) \equiv f_{n}(\bm{a}) + s\gamma\nabla_1.\pmod{\alpha} \]
\end{theorem}
\begin{proof}
    From \cref{cor:remainder-after-change-reduce-vector-modulo-alpha} and \cref{prop:nabla-are-equivalent-modulo-alpha}, we have
    \begin{align*}
        f_{n}(\bm{a} + s\Phi_k\bm{e}_k) &\equiv f_{n}(\bm{a}) + \frac{\gamma}{\alpha^{k}}\cdot(\beta^{s\Phi_k} - 1) &&\pmod{\alpha}\\
        &\equiv f_{n}(\bm{a}) + \gamma\cdot\nabla_{k}\cdot\sum_{i=0}^{s-1}\beta^{i\Phi_k} &&\pmod{\alpha} \\
        &\equiv f_{n}(\bm{a}) + \gamma\cdot\nabla_1\cdot\sum_{i=0}^{s-1} 1^{i} &&\pmod{\alpha} \\
        &\equiv f_{n}(\bm{a}) + s\gamma\nabla_1. &&\pmod{\alpha}
    \end{align*}
\end{proof}

Back to \cref{thm:the-periodic-rule-of-reduce-sequence}. In fact, it presents the variaty of $s$ that preserve integers, which do not depend on the original reduced sequence $\bm{v}$, but only on the position $k$.
Therefore, we can classify these reduced sequences into different equivalence classes.
Furthermore, this conclusion can be easily generalized to the study of the periodic properties of linear or fractional dynamical systems in other algebraic fields, simply by replacing the prime factors with ideals and the LCM with the least common ideal, while the form remains essentially unchanged.
However, such generalizations are beyond the scope of this paper, which will continue to focus on Collatz sequences.
